% Session 3 -- Preparing the Run: Procedure TPS-P-02, the parameter record
% (Operating Procedure OP-7.3, which messages TPS-001 and TPS-002 name).
% Law 1 sources: spec/tps_ritual.json R1 (TPSPARM.DAT opened through FCB
% AH=0Fh in the current directory; TITLE=, DIVISION=, MEMO=, LINES=; refuse
% TPS-001 / refuse_missing_key TPS-002; exit 1); os/milton/tpsrep_program.asm
% (R1 before the table is opened; getparam asks TITLE, DIVISION, MEMO,
% LINES in that order, so the first missing name is the one reported; the
% banner first); os/milton/command.c (COPY CON, Ctrl-Z ends the file).
% Screens tps001, tps002, tpsstart: the journey route.
% GATES: test-tps-report test-flair-tps-journey test-pubs-screens test-pubs-tps
\Gates{test-tps-report, test-flair-tps-journey, test-pubs-screens, test-pubs-tps}
\chapter{Preparing the Run}

\SideHead{Introduction}The first thing the TPS Report Generator does is
read its parameter record, \file{TPSPARM.DAT}, from the current directory.
A run without a complete parameter record ends at once.  In this session you
start the program and verify that the parameter record has been accepted.
The parameter record itself is maintained under Operating Procedure OP-7.3
(Parameter Files), which is \proc{02}.

\SideHead{Objectives}Upon completion of this session you will be able to:
\begin{enumerate}
\item Start the TPS Report Generator from the prompt \texttt{A:\bs APPS>}.
\item Recognize the messages TPS-001 and TPS-002 and take the action they
  require.
\item State what each line of the parameter record controls.
\end{enumerate}

\section{The Parameter Record}
The four lines of the parameter record are described in Session~1.  The
record on the shipped volume is reproduced in Figure~1-1.  The operator does
not change the parameter record.  When a line must change, for example when
a new memo about the cover sheets is issued, the change is made as your
supervisor directs, and the new record is filed in \file{A:\bs APPS} under
the same name.

The program reads the lines in this order: \texttt{TITLE=},
\texttt{DIVISION=}, \texttt{MEMO=}, \texttt{LINES=}.  The first line that is
missing is the one named in the message.

\begin{Procedure}{TPS-P-02}{To Start the Program (OP-7.3)}
\ProcField{Purpose}{To start the TPS Report Generator and confirm that it has
  accepted the parameter record.}
\ProcField{Frequency}{Each run, directly after \proc{01}.}
\ProcField{Performed by}{The operator.  Changes to the parameter record: as
  your supervisor directs.}
\ProcField{Prerequisites}{\proc{01} is complete; the prompt is
  \texttt{A:\bs APPS>}.}
\ProcField{Materials}{\form{5}.}
\begin{ProcSteps}
\item Type \type{TPSREP TPS.DBF} and press \key{Enter}.  The program
  displays its banner.
  \OpsShot[fromtext={INITECH CORPORATION}, totext={(C) 1989}]{tpsstart}
\item Read the next line of the screen.
  \begin{itemize}
  \item If it is the first line of the form alignment prompt, the parameter
    record has been accepted.  Continue with \proc{03}.
  \item If it is a message beginning \texttt{TPS-001} or \texttt{TPS-002},
    the run has ended with return code 1.  Go to step 3.
  \end{itemize}
\item Enter the message number, the date and the time on \form{5}, line 2,
  and see your supervisor.  Do not continue the run.
\end{ProcSteps}
\ProcField{Verification}{The banner is followed by the form alignment
  prompt.}
\ProcField{Records retained}{\form{5}, line 2 (a refusal only).}
\ProcNote{The banner is displayed whether or not the parameter record is
  accepted.  The line after the banner decides.}
\end{Procedure}

\section{When the Record Is Not Found}
\index{TPS-001}
If there is no parameter record in the current directory, the program
displays:
\TpsMsg{TPS-001}
and ends with return code 1.  This is what happens when the program is run
from a directory other than \file{A:\bs APPS}:

\OpsShot[fromtext={A:\Bsl>APPS}, label=fig:tps001, caption={The program run
from the root directory.}]{tps001}

\SideHead{Action}Make \file{A:\bs APPS} the current directory (\proc{01},
step 6) and start again at \proc{02}, step 1.  If the message appears at the
prompt \texttt{A:\bs APPS>}, see your supervisor.

\section{When a Line Is Missing}
\index{TPS-002}
If the parameter record lacks one of its four lines, the program displays:
\TpsMsg{TPS-002}
where \ph{n} is the name of the missing line, and ends with return code 1.
In the example below a record of two lines has been made in the root
directory, and the program run from there, to show the message.  The line
named is \texttt{MEMO}, the first of the missing lines in reading order:

\OpsShot[fromtext={A:\Bsl>COPY CON}, caption={A parameter record without its
\texttt{MEMO=} and \texttt{LINES=} lines.}]{tps002}

\SideHead{Action}See your supervisor.  The record is corrected under OP-7.3.

\section{Lab Project 2. Starting the Program}
During this project you start the program in the wrong directory and then
in the right one.
\begin{enumerate}
\item Carry out \proc{01}, steps 1 to 5, so that the prompt is
  \texttt{A:\bs>}.
\item Type \type{APPS\bs TPSREP TPS.DBF} and press \key{Enter}.  Compare
  your screen with Figure~\ref{fig:tps001}.  Note the return code:
  \rule{20mm}{0.4pt}
\item Type \type{CD APPS} and press \key{Enter}.
\item Carry out \proc{02}, step 1.  The form alignment prompt follows the
  banner.  Continue with Lab Project 3 (Session~4) without leaving the
  program.
\end{enumerate}

\section{Study Questions}
\begin{enumerate}
\item From which directory does the program read its parameter record?
\item What message is displayed when the parameter record is not found, and
  what return code does the program end with?
\item A parameter record has only its \texttt{TITLE=} and \texttt{LINES=}
  lines.  Which line does message TPS-002 name?
\item Who directs a change to the parameter record?
\item What does the line after the banner tell you?
\end{enumerate}
